Многое, особенно сведения из разделов про тензорные произведения, взято из книги~\cite{gorodentsev03}. Очень рекомендую её и предшествующую ей книгу~\cite{gorodentsev01}. Они замечательные, как и весь курс А.\,Л.~Городенцева по алгебре, прочитанный в НМУ в 2020--2022 годах.

\paragraph{Список обозначений.}
\begin{itemize}
	\item $[n] := \set*{1, 2, \dots, n}$, если только квадратные скобки не означают класс по модулю вычетов; что имеется в виду, обычно ясно из контекста;
	\item если $I := (i_1 < i_2 < \dots < i_k) \subset [n]$ --- упорядоченный набор, то
		$
			\hat{I} := [n] \setminus I
		$
		состоит из тех чисел из~$[n]$, которые не вошли в~$I$, и упорядочен по возрастанию;
	\item если $i \in [n]$, то вместо $[n] \setminus \set*{i}$ будем писать $[n] \setminus i$.
\end{itemize}
